\section{Generalization Beyond Classical Expectations}

The most famous theoretical puzzle of deep learning concerns generalization.
Why do large networks often predict well on new data despite extreme flexibility?

\subsection*{Fitting versus predicting}

The training objective is usually based on empirical risk:
\[
\widehat{R}_n(f)
=
\frac{1}{n}\sum_{i=1}^n \ell(f(x_i),y_i).
\]
The quantity we ultimately care about is the population risk:
\[
R(f)=\mathbb{E}_{(x,y)\sim P}[\ell(f(x),y)].
\]

\medskip

Classical intuition emphasizes the danger that low empirical risk may not imply low population risk when the model class is too rich.
That danger is real.
But in modern deep learning, interpolation of the training data is often compatible with surprisingly good test performance.

\subsection*{Implicit bias of optimization}

One important idea is that learning algorithms do not explore parameter space arbitrarily.
Gradient-based optimization may prefer certain kinds of solutions over others, even when many parameter settings achieve nearly identical training error.

\medskip

This preference is often called an \emph{implicit bias} or \emph{implicit regularization} of optimization.

\medskip

Thus the learned predictor is not determined only by:
\begin{itemize}
    \item the model class,
    \item the dataset,
    \item the explicit loss.
\end{itemize}
It is also shaped by:
\begin{itemize}
    \item initialization,
    \item optimization dynamics,
    \item architecture,
    \item training schedule,
    \item normalization and regularization mechanisms.
\end{itemize}

This helps explain why overparameterized systems may still converge to useful solutions rather than arbitrary memorizing ones.

\subsection*{Norm and margin viewpoints}

Another family of explanations emphasizes geometric quantities such as norm and margin.

\paragraph{Norm.}
A predictor may interpolate the training data while still being ``simple'' in an appropriate geometric sense, for example by having controlled parameter norm or controlled function complexity.

\paragraph{Margin.}
In classification, correct predictions made with large margin are often more stable than predictions made barely on the correct side of a boundary.
Margin-based reasoning was central in classical methods such as support vector machines, and related ideas continue to matter in deep learning.

\medskip

These viewpoints do not provide a complete explanation of generalization in all deep models, but they often provide useful partial structure.

\subsection*{Overparameterization as a changed regime}

A striking modern lesson is that increasing model size does not always hurt test performance.
In some settings, once the model becomes large enough to fit the data well, further increases in size may actually improve generalization again.

\medskip

This suggests that the classical ``more complex model means worse generalization'' picture is too simplistic for highly overparameterized regimes.

\subsection*{Double descent}

One influential conceptual picture is \emph{double descent}.
In a simplified description, as model size increases:
\begin{itemize}
    \item test error may first decrease,
    \item then increase near an interpolation threshold,
    \item then decrease again in a larger overparameterized regime.
\end{itemize}

This phenomenon does not occur identically in all settings, but it captures a real and important departure from classical bias--variance intuition.

\subsection*{What appears established}

Several broad points now seem plausible and robust:
\begin{itemize}
    \item interpolation alone does not determine generalization quality,
    \item optimization dynamics matter,
    \item geometric notions such as norm and margin often remain relevant,
    \item overparameterization changes the effective learning regime rather than simply making it worse.
\end{itemize}

\subsection*{What remains unclear}

At the same time, no single theory fully explains:
\begin{itemize}
    \item why stochastic gradient methods often find good solutions so reliably,
    \item why overparameterized models can be simultaneously easy to optimize and strong at prediction,
    \item what depth contributes beyond width or linearized approximations,
    \item how architecture and data structure enter the generalization story in a unified way.
\end{itemize}

Thus deep-learning generalization is not mysterious in the sense of being completely opaque, but it remains incomplete in the sense of lacking a definitive unified account.

